\chapter{Hasil dan Pembahasan}

Bab ini menyajikan hasil rancang bangun sistem Rupiah Digital ritel berbasis
Hyperledger Fabric beserta pembahasannya. Bagian pertama memaparkan hasil
implementasi sistem yang menjawab tujuan pertama dan kedua, bagian kedua
memaparkan hasil evaluasi kinerja yang menjawab tujuan ketiga, dan bagian
ketiga membandingkan hasil penelitian dengan penelitian terdahulu.

\section{Hasil Implementasi Sistem}

\subsection{Jaringan Hyperledger Fabric}

Jaringan \textit{blockchain} berhasil dibangun dengan lima organisasi peserta
(BankIndonesiaOrg, HimbaraBankOrg, CommercialBankOrg, PJPOrg, dan
OJKObserverOrg) serta satu organisasi \textit{orderer}. Setiap organisasi
peserta menjalankan satu \textit{peer} dengan basis data status CouchDB, dan
layanan pengurutan menggunakan konsensus Raft (etcdraft). Seluruh komponen
dijalankan sebagai \textit{container} melalui Docker Compose. Material
kriptografi, blok genesis, dan konfigurasi kanal dibangkitkan secara otomatis
melalui skrip, kemudian \textit{chaincode} \texttt{digital-rupiah} dipaketkan,
dipasang pada seluruh organisasi, disetujui sesuai kebijakan dukungan
\textit{majority}, dan di-\textit{commit} ke kanal. Dengan kebijakan dukungan
\textit{majority}, setiap pemutakhiran \textit{ledger} memerlukan dukungan
mayoritas organisasi sehingga tidak ada satu pihak yang dapat mengubah status
secara sepihak.

\subsection{Chaincode Digital Rupiah}

\textit{Chaincode} \texttt{digital-rupiah} mengimplementasikan logika bisnis
Rupiah Digital pada bahasa Go. Fungsi-fungsi \textit{chaincode} dikelompokkan
menurut domainnya sebagaimana dirangkum pada Tabel~\ref{tab:chaincode}. Seluruh
operasi yang mengubah status \textit{ledger} mencatatkan jejak audit dan/atau
\textit{event} pengawasan sehingga dapat ditelusuri oleh pengawas.

\begingroup
\renewcommand{\arraystretch}{1.3}
\begin{longtable}{p{3cm} p{10cm}}
\caption{Kelompok fungsi chaincode digital-rupiah}\label{tab:chaincode}\\
\toprule
\textbf{Kelompok} & \textbf{Fungsi utama} \\
\midrule
\endfirsthead
\toprule
\textbf{Kelompok} & \textbf{Fungsi utama} \\
\midrule
\endhead
\bottomrule
\endlastfoot
Peserta & \texttt{SubmitParticipant}, \texttt{ApproveParticipant},
\texttt{FreezeParticipant}, \texttt{UnfreezeParticipant},
\texttt{RejectParticipant}, \texttt{OffboardParticipant} \\
Dompet & \texttt{CreateWallet}, \texttt{CreateWholesaleWallet}, \texttt{Mint},
\texttt{Burn}, \texttt{Transfer}, \texttt{FreezeWallet},
\texttt{UnfreezeWallet} \\
Likuiditas & \texttt{RequestIssuance}, \texttt{RequestIssuanceRtgs},
\texttt{DistributeToParticipant}, \texttt{RequestRedemption} \\
Limit & \texttt{SetTierLimit}, \texttt{SetSystemLimit},
\texttt{SetAutoLimitPolicy} \\
KYC & \texttt{SubmitKycProfile}, \texttt{RefreshKycProfile},
\texttt{GetKycProfile} \\
Pengawasan & \texttt{GetReconciliationReport}, \texttt{GetMetrics},
\texttt{GetSupervisionEvents}, \texttt{GetTransactionHistory} \\
\end{longtable}
\endgroup

Penegakan kepatuhan dilakukan langsung pada \textit{chaincode}. Tingkatan
dompet ritel diturunkan dari jenis subjek, tingkat uji tuntas, tingkat risiko,
dan persetujuan senior pada profil KYC. Operator tidak dapat memilih tingkatan.
Dompet institusional tidak menggunakan tingkatan ritel dan tunduk pada batas
tingkat sistem. Tabel~\ref{tab:tier} menunjukkan parameter kebijakan purwarupa,
sedangkan Gambar~\ref{fig:participant-lifecycle} menunjukkan siklus hidup
peserta dari pendaftaran hingga pelepasan keanggotaan.

\begin{figure}[H]
\centering
\resizebox{0.88\textwidth}{!}{%
\begin{tikzpicture}[
  shorten >=1pt, node distance=3.6cm, on grid, auto,
  state/.style={draw=black, circle, fill=white, minimum size=1.8cm,
    align=center, font=\footnotesize\bfseries},
  ar/.style={-{Stealth}, thick},
  lab/.style={font=\scriptsize, align=center},
]
\node[state] (pending)    at (0, 0)    {pending};
\node[state] (active)     at (4.0, 0)  {active};
\node[state] (frozen)     at (8.0, 0)  {frozen};
\node[draw=black, dashed, rounded corners=3pt, fill=gray!12, align=center,
      minimum height=1.5cm, minimum width=2.8cm, font=\scriptsize] (deleted) at (2.0,-3.0)
      {Dihapus dari\\\textit{ledger}\\(\texttt{DelState})};

\draw[ar] (pending.east)  -- node[above, lab]{ApproveParticipant} (active.west);
\draw[ar] (active.east)   -- node[above, lab]{FreezeParticipant} (frozen.west);
\draw[ar] (frozen.south west) to[out=-150, in=30]
  node[below, lab]{UnfreezeParticipant} (active.south east);
\draw[ar] (active.south)  to[out=-90, in=20]
  node[right, lab]{OffboardParticipant} (deleted.north east);
\draw[ar] (pending.south) to[out=-90, in=160]
  node[left, lab]{RejectParticipant} (deleted.north west);
\end{tikzpicture}%
}
\caption{Siklus hidup peserta pada \textit{chaincode}: dari \textit{pending} (menunggu persetujuan) menjadi \textit{active}, kemudian dapat dibekukan (\textit{frozen}) dan dicairkan kembali. Penolakan (\texttt{RejectParticipant}) maupun pelepasan keanggotaan (\texttt{OffboardParticipant}) menghapus rekaman peserta dari \textit{ledger} (\texttt{DelState}), bukan menyimpannya sebagai status}
\label{fig:participant-lifecycle}
\end{figure} Pada fungsi \texttt{Transfer}, sistem memeriksa status
beku dompet, kelayakan KYC pengirim, batas per transaksi, pengeluaran harian,
pengeluaran bulanan, penerimaan bulanan, kecukupan saldo pengirim, serta saldo
maksimum penerima sebelum transaksi dieksekusi.

\begingroup
\footnotesize
\renewcommand{\arraystretch}{1.3}
\setlength{\tabcolsep}{3pt}
\begin{table}[H]
\caption{Batas tingkatan dompet ritel (dalam Rupiah)}
\label{tab:tier}
\centering
\begin{tabular}{@{}l r r r r r@{}}
\toprule
\textbf{Tingkatan} & \textbf{Saldo maks.} & \textbf{Per transaksi} &
\textbf{Harian keluar} & \textbf{Bulanan keluar} & \textbf{Bulanan masuk} \\
\midrule
BASIC & 2.000.000 & 250.000 & 500.000 & 5.000.000 & 20.000.000 \\
STANDARD & 20.000.000 & 2.500.000 & 10.000.000 & 40.000.000 & 40.000.000 \\
MERCHANT & 200.000.000 & 10.000.000 & 50.000.000 & 500.000.000 & 500.000.000 \\
\bottomrule
\end{tabular}
\end{table}
\endgroup

Selain limit per tingkatan, sistem menetapkan batas tingkat sistem berupa
pasokan global maksimum Rp10.000.000.000.000, saldo maksimum per peserta
Rp1.000.000.000.000, dan nilai maksimum per transaksi Rp100.000.000.000.

\subsection{Backend API}

Lapisan \textit{backend} menyediakan antarmuka REST yang menerjemahkan
permintaan HTTP menjadi pemanggilan \textit{chaincode} melalui \textit{Fabric
Gateway SDK}. Kontrol akses berbasis peran ditegakkan melalui \textit{middleware}
dengan tujuh peran sebagaimana dirangkum pada Tabel~\ref{tab:peran}. Pemisahan
peran ini memastikan bahwa, misalnya, hanya peran \texttt{bank\_indonesia} yang
dapat melakukan penerbitan, sedangkan peran
\texttt{supervisor} hanya memiliki akses baca terhadap laporan pengawasan.

\begingroup
\renewcommand{\arraystretch}{1.3}
\begin{table}[H]
\caption{Peran akses pada backend API}
\label{tab:peran}
\centering
\begin{tabular}{p{3.2cm} p{9.8cm}}
\toprule
\textbf{Peran} & \textbf{Kewenangan utama} \\
\midrule
\texttt{public} & Pemeriksaan kesehatan dan topologi jaringan \\
\texttt{authen\allowbreak ticated} & Pembacaan daftar dompet \\
\texttt{kyc\_verified} & Transfer dan pembacaan saldo pelanggan yang telah disetujui \\
\texttt{bank\_pjp} & Onboarding pelanggan, keputusan KYC, pembuatan dompet, dan operasi peserta \\
\texttt{bank\allowbreak\_indonesia} & Penerbitan, kebijakan limit, dan pengawasan; akses baca KYC \\
\texttt{merchant} & Transfer dan pembacaan saldo pedagang yang telah disetujui \\
\texttt{supervisor} & Akses baca KYC, laporan, dan pengawasan \\
\bottomrule
\end{tabular}
\end{table}
\endgroup

\subsection{Hasil Pengujian Fungsional}

Pengujian fungsional dilakukan untuk memverifikasi setiap alur proses yang
dirancang pada BAB III. Hasil pengujian menunjukkan bahwa seluruh alur berjalan
sesuai rancangan, yaitu:
\begin{enumerate}
    \item \textbf{Penerbitan dan distribusi berjenjang}. Bank Indonesia
    menerbitkan Rupiah Digital ke peserta; bank validator mendistribusikan
    likuiditas ke PJP melalui \texttt{DistributeToParticipant} yang melakukan
    \texttt{Burn} pada pengirim dan \texttt{Mint} pada penerima; PJP kemudian
    menyalurkan saldo ke pelanggan ritel. Upaya penerbitan langsung kepada PJP
    ditolak sesuai aturan, sehingga PJP hanya memperoleh likuiditas melalui
    bank.
    \item \textbf{Penegakan limit}. Transaksi yang melampaui batas per
    transaksi, keluar harian, keluar bulanan, atau masuk bulanan ditolak,
    demikian pula transfer yang menyebabkan saldo penerima melampaui batas.
    \item \textbf{KYC dan pembekuan}. Bank/PJP mengajukan dan memperbarui
    profil KYC. BASIC diturunkan untuk pelanggan risiko rendah dengan uji tuntas
    sederhana; STANDARD untuk pelanggan dengan uji tuntas standar/lanjutan;
    MERCHANT untuk pedagang dengan uji tuntas standar/lanjutan. Risiko tinggi
    memerlukan EDD dan persetujuan senior. Dompet yang dibekukan tidak dapat
    bertransaksi.
    \item \textbf{Transfer ritel}. Jalur \texttt{Transfer} yang sama berhasil
    digunakan untuk transaksi antarpelanggan dan pelanggan-ke-pedagang.
    \item \textbf{Pengawasan}. Laporan rekonsiliasi, metrik, dan \textit{event}
    pengawasan dapat diambil oleh peran berwenang.
\end{enumerate}

\section{Hasil Evaluasi Kinerja}

Evaluasi kinerja menggunakan Hyperledger Caliper 0.6 pada jaringan lokal lima
organisasi dengan sapuan konfigurasi \textit{worker} (1, 2, dan 4
\textit{worker}) untuk mengidentifikasi batas pemrosesan sisi server. Setiap
profil transfer menjalankan 40 transaksi pada laju kirim sekitar 10 TPS
(pemanasan), 30 TPS (stabil), dan 60 TPS (puncak). Dua puluh persen pelanggan
memperoleh BASIC melalui uji tuntas sederhana, sedangkan 80\% memperoleh
STANDARD melalui uji tuntas standar; pedagang memperoleh MERCHANT.

Beban terukur hanya memakai jalur yang masih berada dalam lingkup penelitian:
70\% transfer antarpelanggan, 25\% transfer pelanggan-ke-pedagang, dan 5\%
pembacaan dompet. Setiap subjek diberi jangkar KYC yang disetujui sebelum ronde
dimulai. Dana awal STANDARD sebesar Rp5.000.000 dan BASIC sebesar Rp500.000
menyisakan ruang penerimaan di bawah saldo maksimum. Pengukuran melaporkan jumlah
berhasil dan gagal, \textit{throughput}, serta latensi minimum, rata-rata, dan
maksimum. Nomor slot global menentukan jenis transaksi secara deterministik dan
setiap kunci dompet yang dimutasi hanya digunakan satu kali dalam satu ronde
terukur. Kegagalan kebijakan dan konflik MVCC tetap dihitung jika terjadi.

Hasil diperoleh dari laporan Caliper setelah \textit{chaincode} versi 2.0
dipasang pada \textit{ledger} bersih dan \texttt{InitLedger} dijalankan. Angka
hasil tidak diwariskan dari beban kerja versi terdahulu karena semantik transaksi,
batas dompet, dan strategi pemilihan kunci berbeda. Prosedur ini mencegah
perbandingan dua beban kerja yang tidak setara dan menjaga keterlacakan hasil.

Konfigurasi pada Tabel~\ref{tab:laju} merupakan beban ritel yang menyerupai
perilaku sistem CBDC ritel pada tingkat purwarupa, bukan klaim lalu lintas
produksi nasional. Kesesuaiannya berasal dari tiga aspek. Pertama, transaksi
yang diukur adalah transfer antarpelanggan, transfer pelanggan-ke-pedagang,
pembacaan dompet, KYC, operasi moneter, kueri pengawasan, dan kebijakan
administratif, sesuai fungsi pembayaran umum pada CBDC ritel
\cite{Auer2020The}. Kedua, populasi pengujian memuat pelanggan BASIC dan
STANDARD serta pedagang MERCHANT dengan KYC yang disetujui, mengikuti prinsip
dompet bertingkat dan limit berbasis identifikasi yang juga muncul pada e-CNY
dan eNaira \cite{pboc2021ecny, cbn2021enaira}. Ketiga, laju tetap 10, 30, dan
60 TPS pada profil transfer memungkinkan observasi hubungan antara
\textit{send rate}, \textit{throughput}, latensi, dan kegagalan transaksi
sebagaimana dianjurkan dalam metrik kinerja blockchain
\cite{hlf_metrics, hlf_caliper}. Dengan batasan tersebut, hasil Caliper layak
digunakan untuk menilai kelayakan desain dan membangun garis dasar pengujian
lanjutan, tetapi tidak untuk mengekstrapolasi kapasitas nasional tanpa data
transaksi riil, topologi produksi, dan perangkat keras setara.

\begin{table}[H]
\caption{Konfigurasi evaluasi transfer ritel}
\label{tab:laju}
\centering
\begin{tabular}{l r r r l}
\toprule
\textbf{Profil} & \textbf{\textit{Worker}} & \textbf{Tx} & \textbf{Laju} & \textbf{Komposisi} \\
\midrule
Transfer & 1/2/4 & 40 & 10 TPS & 70\% P2P, 25\% pelanggan--pedagang, 5\% baca \\
Transfer & 1/2/4 & 40 & 30 TPS & 70\% P2P, 25\% pelanggan--pedagang, 5\% baca \\
Transfer & 1/2/4 & 40 & 60 TPS & 70\% P2P, 25\% pelanggan--pedagang, 5\% baca \\
Fungsional & 2 & 40--80 & 5--30 TPS & KYC, moneter, pengawasan, administratif \\
\bottomrule
\end{tabular}
\end{table}

Pengujian dijalankan pada jaringan Garuda yang diatur ulang, \textit{chaincode}
\texttt{digital-rupiah} versi 2.0 dideploy ulang dengan
\texttt{NETWORK\_MODE=garuda}, kemudian \texttt{InitLedger} dijalankan dan
kontainer \textit{chaincode} dipanaskan sebelum Caliper dieksekusi. Laporan
pengujian menggunakan cap waktu \texttt{2026-06-30-113543}. Seluruh 616
permintaan \textit{chaincode} yang dilaporkan Caliper berhasil dengan 0 kegagalan.
Log \textit{peer} tidak menunjukkan konflik MVCC maupun penolakan kebijakan/limit.

\begin{table}[H]
\caption{Hasil evaluasi transfer ritel Hyperledger Caliper}
\label{tab:hasil-kinerja-caliper}
\centering
\scriptsize
\begin{tabular}{l r r r r r r r r}
\toprule
\textbf{Ronde} & \textbf{\textit{Worker}} & \textbf{Sukses} & \textbf{Gagal} & \textbf{Kirim} &
\textbf{\textit{Throughput}} & \multicolumn{3}{c}{\textbf{Latensi (detik)}} \\
\cmidrule(lr){7-9}
 &  &  &  & \textbf{(TPS)} & \textbf{(TPS)} & \textbf{Min.} &
\textbf{Rerata} & \textbf{Maks.} \\
\midrule
Pemanasan & 1 & 40 & 0 & 10,3 & 7,8 & 0,01 & 0,78 & 2,14 \\
Stabil & 1 & 40 & 0 & 30,8 & 12,7 & 0,00 & 0,62 & 2,15 \\
Puncak & 1 & 40 & 0 & 57,6 & 14,5 & 0,07 & 0,66 & 2,32 \\
Pemanasan & 2 & 40 & 0 & 10,5 & 7,8 & 0,01 & 0,80 & 2,14 \\
Stabil & 2 & 40 & 0 & 31,6 & 12,3 & 0,01 & 0,64 & 2,25 \\
Puncak & 2 & 40 & 0 & 61,0 & 14,5 & 0,03 & 0,69 & 2,27 \\
Pemanasan & 4 & 40 & 0 & 11,1 & 8,0 & 0,01 & 0,81 & 2,18 \\
Stabil & 4 & 40 & 0 & 33,3 & 12,9 & 0,01 & 0,64 & 2,16 \\
Puncak & 4 & 40 & 0 & 65,9 & 14,5 & 0,01 & 0,70 & 2,30 \\
\bottomrule
\end{tabular}
\end{table}

Pada profil transfer, seluruh kombinasi 1, 2, dan 4 \textit{worker} menghasilkan
40 permintaan sukses dan 0 kegagalan pada setiap ronde. Peningkatan laju kirim
dari sekitar 10 TPS ke 60 TPS tidak menaikkan \textit{throughput} secara linear.
\textit{Throughput} puncak berada pada sekitar 14,5 TPS untuk ketiga jumlah
\textit{worker}. Hal ini menunjukkan bahwa pada topologi lokal, batas pemrosesan
lebih dipengaruhi oleh jaringan Fabric dan proses validasi dibanding jumlah
\textit{worker} Caliper. Keberhasilan 100\% pada pengujian berdurasi terbatas ini
menunjukkan bahwa transaksi valid dapat diselesaikan tanpa konflik kunci pada
konfigurasi yang dilaporkan, bukan bahwa sistem siap menangani beban nasional.

\begin{table}[H]
\caption{Hasil evaluasi fungsi lain Hyperledger Caliper}
\label{tab:hasil-kinerja-fungsional-caliper}
\centering
\scriptsize
\begin{tabular}{l r r r r r r r}
\toprule
\textbf{Profil} & \textbf{Sukses} & \textbf{Gagal} & \textbf{Kirim} &
\textbf{\textit{Throughput}} & \multicolumn{3}{c}{\textbf{Latensi (detik)}} \\
\cmidrule(lr){6-8}
 &  &  & \textbf{(TPS)} & \textbf{(TPS)} & \textbf{Min.} &
\textbf{Rerata} & \textbf{Maks.} \\
\midrule
KYC/\textit{onboarding} & 56 & 0 & 4,4 & 3,7 & 0,13 & 0,88 & 2,29 \\
Operasi moneter & 80 & 0 & 20,5 & 19,9 & 0,10 & 0,32 & 0,53 \\
Kueri pengawasan & 80 & 0 & 30,8 & 17,0 & 0,01 & 1,10 & 3,24 \\
Kebijakan administratif & 40 & 0 & 5,3 & 5,2 & 0,09 & 0,91 & 1,72 \\
\bottomrule
\end{tabular}
\end{table}

Tabel~\ref{tab:hasil-kinerja-fungsional-caliper} menunjukkan bahwa fungsi
nontransfer juga dapat dieksekusi tanpa kegagalan pada beban terbatas. Profil
KYC/\textit{onboarding} berisi beberapa pemanggilan \textit{chaincode} untuk satu
alur bisnis, sehingga angka sukses pada tabel dibaca sebagai jumlah permintaan
\textit{chaincode} sukses, bukan jumlah lengkap pelanggan yang di-\textit{onboard}.
Kueri pengawasan memiliki \textit{throughput} lebih rendah daripada laju kirim
karena sebagian kueri membaca himpunan data buku besar, sedangkan operasi moneter
memiliki latensi rata-rata paling rendah pada pengujian ini.

\section{Perbandingan Hasil Penelitian dengan Hasil Terdahulu}

Perbandingan dengan purwarupa token sederhana berbasis Hyperledger Fabric milik
Hanif \cite{hanif2024meetcoin} dirangkum pada
Tabel~\ref{tab:perbandingan-hanif}. Data penelitian terdahulu dipakai untuk
menjelaskan perkembangan rancangan, bukan sebagai klaim kapasitas yang setara.

Penelitian ini menambahkan distribusi dua jenjang, pemisahan data KYC luring dan
jangkar kepatuhan pada buku besar, penurunan tingkatan dompet berbasis risiko,
limit masuk dan keluar, serta peristiwa pengawasan. Perbandingan angka kinerja
langsung tidak dilakukan karena perangkat keras, topologi, kebijakan dukungan,
dan semantik beban kerja berbeda.

\begin{table}[H]
    \caption{Perbandingan purwarupa MeetCoin dan purwarupa Rupiah Digital ritel. Sumber penelitian terdahulu: Hanif \cite{hanif2024meetcoin}.}
    \label{tab:perbandingan-hanif}
    \centering
    \scriptsize
    \begin{tabular}{@{}p{2.5cm}p{4.6cm}p{5.1cm}@{}}
        \toprule
        \textbf{Dimensi} & \textbf{MeetCoin} & \textbf{Penelitian ini} \\
        \midrule
        Domain & Token dan dompet digital sederhana & Purwarupa CBDC ritel dua jenjang \\
        Topologi & Dua \textit{peer}, satu \textit{orderer}, dan satu kanal & Lima organisasi peserta, satu organisasi \textit{orderer}, dan satu kanal \\
        Basis data status & LevelDB berbasis kunci--nilai & CouchDB berbasis dokumen JSON dan kueri kaya \\
        Identitas dan kepatuhan & Fabric CA dan autentikasi JWT tanpa tingkatan KYC & KYC luring, jangkar kepatuhan, serta tingkatan BASIC, STANDARD, dan MERCHANT \\
        Beban kerja & \textit{Mint}, cek saldo, dan transfer; 5/10/15 \textit{worker}; 1.000/2.000/3.000 transaksi; 50/100/150 TPS & 70\% transfer antarpelanggan, 25\% pelanggan--pedagang, dan 5\% baca dompet dengan dua identitas organisasi \\
        Metrik & Keberhasilan, \textit{throughput}, latensi, CPU, dan memori & Keberhasilan, \textit{throughput}, latensi, serta konflik MVCC \\
        \bottomrule
    \end{tabular}
\end{table}

\paragraph{Perbandingan dengan PoC Proyek Garuda.}
Laporan PoC Proyek Garuda \cite{bi_garuda_poc} mendokumentasikan uji coba
\textit{wholesale Rupiah Digital} (wRD) pada tahap \textit{Immediate State}.
Berbeda dengan penelitian ini, PoC Garuda bersifat grosir dan mengevaluasi dua
platform non-Fabric: R3 Corda (UTXO, konsensus Notary) dan Kaleido Hyperledger
Besu (akun, konsensus QBFT/PoA), yang dipilih dari 39 kandidat menjadi dua
finalis. Pengujian dilakukan dengan 55 skenario bisnis dan teknis menggunakan
Locust Dashboard serta AWS CloudWatch. Hasil \textit{load test} melaporkan
\textit{throughput} 37 TPS (Corda) dan 39 TPS (Besu), keduanya memenuhi ambang
batas lalu lintas grosir sekitar 30 TPS. Arsitektur yang diuji terdiri atas enam
lapisan (\textit{use case}, aset digital, eksekusi, data, konsensus, jaringan)
dengan aspek keamanan menyeluruh, serta peran Master Node BI sebagai KDR
(\textit{Khazanah Digital Rupiah}), Regulator, Observer, Administrator, dan
Provider. PoC Garuda tidak mengevaluasi dompet ritel berjenjang, KYC bertingkat,
pembayaran luring, atau kanal pelanggan--pedagang karena lingkupnya terbatas
pada penerbitan, penebusan, dan transfer dana antarlembaga. Penelitian ini
mengisi celah ritel tersebut dengan Hyperledger Fabric, dompet BASIC, STANDARD,
dan MERCHANT, serta pengujian Caliper pada beban transfer ritel. Perbedaan
platform (Fabric vs Corda/Besu), alat ukur (Caliper vs Locust), dan lingkup
(ritel vs grosir) membuat angka TPS kedua penelitian tidak dapat dibandingkan
secara langsung, tetapi saling melengkapi gambaran kelayakan DLT untuk CBDC
Rupiah Digital pada lapis grosir dan ritel. Distribusi dua jenjang yang
diarahkan Bank Indonesia \cite{bi_whitepaper2022, bi_consultative} tetap
dipertahankan dalam purwarupa ini.

\paragraph{Perbandingan dengan tinjauan literatur Aryani (2023).}
Aryani \cite{aryani2023critlit} melakukan tinjauan literatur kritis komparatif
terhadap \textit{White Paper Project Garuda} dan implementasi CBDC di empat
negara (eNaira, Sand Dollar, e-CNY, Project Orchid). Kajian tersebut
merekomendasikan adopsi dompet berjenjang dan KYC bertingkat untuk r-Digital
Rupiah, dengan mengutut eNaira (lima tingkat), Sand Dollar (dua tingkat), dan
e-CNY (berjenjang berdasarkan kekuatan identifikasi). Penelitian ini
mengimplementasikan rekomendasi tersebut sebagai kebijakan yang dapat diuji:
tiga tingkatan dompet (BASIC, STANDARD, MERCHANT) diturunkan dari jenis subjek,
tingkat uji tuntas, tingkat risiko, dan persetujuan senior pada profil KYC.
Aryani juga mengutip metrik kuantitatif dari literatur global yang menunjukkan
rata-rata DLT privat hanya mencapai sekitar 20 TPS, jauh di bawah jaringan
terpusat seperti VISA (hingga 65.000 TPS), serta menyatakan keprihatinan bahwa
skalabilitas DLT mungkin menghambat \textit{settlement} ritel sehingga
r-Digital Rupiah mungkin memerlukan infrastruktur terpusat. Penelitian ini
memberikan bukti empiris pada titik tersebut: pengujian Caliper pada Hyperledger
Fabric menghasilkan \textit{throughput} transfer 7,8--14,5 TPS dan operasi
fungsional 3,7--19,9 TPS, yang konsisten dengan rentang rata-rata DLT privat
yang dikutip Aryani dan menggarisbawahi perlunya evaluasi topologi produksi
sebelum klaim kapasitas nasional. Seluruh 616 permintaan berhasil tanpa konflik
MVCC maupun penolakan kebijakan, meskipun dengan keterbatasan satu inah dan
topologi purwarupa yang belum mewakili produksi nasional. Selain itu, Aryani mencatat bahwa BI Reg.~20/6/PBI/2018
membatasi uang elektronik tak terdaftar hingga Rp2.000.000 dan terdaftar hingga
Rp10.000.000, yang dapat berfungsi sebagai proksi limit dompet berjenjang;
parameter Rupiah pada purwarupa ini dinyatakan sebagai pilihan desain dengan
ketentuan uang elektronik domestik hanya sebagai proksi dan bukan dasar hukum
CBDC \cite{bi2022padg247}.

\paragraph{Perbandingan dengan e-CNY dan eNaira.}
e-CNY dan eNaira menunjukkan bahwa klasifikasi dompet dan limit dapat dikaitkan
dengan kekuatan identifikasi pengguna \cite{pboc2021ecny, cbn2021enaira}.
Penelitian ini mengadopsi prinsip tersebut sebagai kebijakan yang dapat diuji,
tetapi tidak menyalin angka limit negara lain. Parameter Rupiah pada purwarupa
dinyatakan sebagai pilihan desain, dengan ketentuan uang elektronik domestik
hanya sebagai proksi dan bukan dasar hukum CBDC \cite{bi2022padg247}.

\begin{table}[H]
    \caption{Perbandingan penelitian ini dengan tinjauan literatur Aryani (2023) dan PoC Proyek Garuda (2024).}
    \label{tab:perbandingan-cbdc}
    \centering
    \scriptsize
    \begin{tabular}{@{}p{2.8cm}p{3.8cm}p{3.8cm}p{3.8cm}@{}}
        \toprule
        \textbf{Dimensi} & \textbf{Aryani (2023)} & \textbf{PoC Garuda (2024)} & \textbf{Penelitian ini} \\
        \midrule
        Jenis CBDC & Ritel (analisis konseptual) & Grosir (PoC eksperimental) & Ritel (purwarupa) \\
        Metodologi & Tinjauan literatur kritis komparatif & PoC dengan 55 skenario bisnis dan teknis & Implementasi purwarupa + Caliper \\
        Platform DLT & Tidak dievaluasi (kutipan dari kasus negara) & R3 Corda + Hyperledger Besu & Hyperledger Fabric \\
        Dompet berjenjang & Direkomendasikan (eNaira 5 tingkat, Sand Dollar 2, e-CNY berjenjang) & Tidak ada (grosir saja) & BASIC, STANDARD, MERCHANT (3 tingkat) \\
        KYC & Direkomendasikan bertingkat & Tidak dibahas eksplisit & Luring, jangkar kepatuhan, penurunan berbasis risiko \\
        Metrik kinerja & Dikutip dari literatur (DLT $\sim$20 TPS, VISA $\sim$65.000 TPS) & Diukur: Corda 37 TPS, Besu 39 TPS (Locust) & Diukur: transfer 7,8--14,5 TPS; fungsional 3,7--19,9 TPS (Caliper) \\
        Lingkup pengujian & Tidak ada pengujian & Penerbitan, penebusan, transfer antarlembaga & Transfer antarpelanggan, pelanggan--pedagang, baca dompet \\
        Pembayaran luring & Direkomendasikan untuk r-Digital Rupiah & Tidak dibahas & Tidak diuji dalam purwarupa ini \\
        Tokenisasi & Didukung (konseptual) & UTXO (Corda) / ERC-20 + ERC-1400 (Besu) & Transfer berbasis kunci--nilai pada Fabric \\
        \bottomrule
    \end{tabular}
\end{table}

\paragraph{Batas generalisasi.}
Purwarupa berjalan pada satu inah, memakai satu \textit{peer} per organisasi,
dan menyimulasikan keputusan penyedia KYC. Oleh karena itu, hasil fungsional
membuktikan konsistensi implementasi terhadap rancangan, bukan kesiapan produksi.
Angka Caliper yang diperoleh berlaku untuk konfigurasi yang dilaporkan dan tidak
dapat langsung digeneralisasi menjadi kapasitas nasional. Keterbatasan ini
sejalan dengan keprihatinan Aryani \cite{aryani2023critlit} bahwa skalabilitas
DLT pada volume ritel nasional belum terbukti dan r-Digital Rupiah mungkin
memerlukan infrastruktur terpusat jika ambang kapasitas DLT terlampaui.
Sebaliknya, PoC Garuda \cite{bi_garuda_poc} hanya menguji lapis grosir hingga
ambang 30 TPS dan belum mengeksplorasi kemampuan maksimum, on-premise, atau
multi-validator; celah ritel yang tertinggal menjadi ruang yang diisi oleh
penelitian ini, tetapi dengan batas purwarupa yang sama seperti di atas.
